\clearpage
\section{Deferred Technical Details}
\label{appx:tech}

\subsection{Pseudocode}
\label{appx:pseudocode}

We give the protocols in order of composition:
Algorithms~\ref{alg:ip}, \ref{alg:hadamard}, and
\ref{alg:permutation} implement the packed reductions of
Section~\ref{sec:blocks}; Algorithm~\ref{alg:packed-pcs}
authenticates their terminal evaluations as in Section~\ref{sec:pcs};
and Algorithms~\ref{alg:single-instance-packing} and
\ref{alg:multiple-instance-packing} compose these components
for the two settings of Section~\ref{sec:together}.
For the intermediate-field variant of Section~\ref{sec:hadamard},
Algorithm~\ref{alg:hadamard} uses the Lagrange basis over
\(\mathbb K\) and computes its quotient via
Eq.~\eqref{eq:hadamard-small-field-aggregation}.
Tabulation changes this computation, not the messages or checks.

% !TEX root = ../main.tex
\begin{algorithm}[!t]
\caption{Packed inner-product proof}
\label{alg:ip}
\KwIn{
\(N=m\kappa\), where \(m, \kappa \in 2^\mathbb{N}\);
\(c\in\Field\);
\(\bm a,\bm b\in\Field^N\)
}
\KwOut{
Terminal values \(\ext{v}_1,\ext{v}_2\) and MLE points
\(\ext{\bm y}_a, \ext{\bm y}_b \in\ExtensionField^{\log_2 N}\);
acceptance bit \(\mathsf{ok}\)
}

Arrange \(\bm a,\bm b\) columnwise as
\(\bm A\leftarrow\concat{\bm a^{(1)},\ldots,\bm a^{(\kappa)}}\)
and \(\bm B\leftarrow\concat{\bm b^{(1)},\ldots,\bm b^{(\kappa)}}
\in\Field^{m\times\kappa}\)\;

\tcp{Compute the auxiliary coefficients over the base field}
$\Prover$ sets:\\
\Indp
\(
    \bm g_A(X)
    \leftarrow
    \sum_{j=1}^{\kappa}
    \bm a^{(j)}X^{j-1}
\);~
\(
    \bm g_B(X)
    \leftarrow
    \sum_{j=1}^{\kappa}
    \bm b^{(j)}X^{j-1}
\)\;

\ForEach{
\(\tau\in
\{-(\kappa-1),\ldots,-1,1,\ldots,\kappa-1\}\)
}{
    Set \(
        f_\tau
        \leftarrow
        \sum_{\substack{j,j'\in[\kappa]\\j-j'=\tau}}
        \left\langle
            \bm a^{(j)},\bm b^{(j')}
        \right\rangle
    \)\;
}

Set
\(
    h_+(X)
    \leftarrow
    \sum_{\tau=1}^{\kappa-1}
    f_\tau X^{\tau-1}
\);~
\(
    h_-(X)
    \leftarrow
    \sum_{\tau=1}^{\kappa-1}
    f_{-\tau}X^{\tau-1}
\)\;
\Indm

\(\PtoV\):
the coefficients of \(h_+(X)\) and \(h_-(X)\)\;

\(\VtoP\):
\(\ext{x}\ChooseRandom\ExtensionField^\times\)\;

\tcp{Reduce to a length-\(m\) extension-field inner product}
$\Prover$ sets
\(
    \ext{\bm a}
    \leftarrow
    \bm g_A(\ext{x})
    \in\ExtensionField^m
\),
\(
    \ext{\bm b}
    \leftarrow
    \bm g_B(\ext{x}^{-1})
    \in\ExtensionField^m
\)\;

$\Verifier$ sets
\(
    \ext{c}
    \leftarrow
    c
    +\ext{x}h_+(\ext{x})
    +\ext{x}^{-1}h_-(\ext{x}^{-1})
\)\;

$\Prover$ and $\Verifier$ apply Fact~\ref{fact:ip} to
\(
    \ext{c}
    \EqOrNot
    \left\langle
        \ext{\bm a},\ext{\bm b}
    \right\rangle,
\)
and obtain terminal values
\(
    \ext{v}_1,\ext{v}_2
\),
the challenge point
\(
    \ext{\bm y}'\in
    \ExtensionField^{\log_2m}
\),
and the acceptance bit \(\mathsf{ok}\)\;

\(\Prover\) and \(\Verifier\) set:\\
\Indp
\(
    \ext{\bm y}_a
    \leftarrow
    \bigl(
        \ext{x}^{2^{\log_2\kappa-1}},
        \ext{x}^{2^{\log_2\kappa-2}},
        \ldots,
        \ext{x}
    \bigr)
    \mathbin{\|}
    \ext{\bm y}' \in \ExtensionField^{\log_2 N}
\)\;
\(
    \ext{\bm y}_b
    \leftarrow
    \bigl(
        \ext{x}^{-2^{\log_2\kappa-1}},
        \ext{x}^{-2^{\log_2\kappa-2}},
        \ldots,
        \ext{x}^{-1}
    \bigr)
    \mathbin{\|}
    \ext{\bm y}' \in \ExtensionField^{\log_2 N}
\)\;
\Indm

Output the terminal claims
\(
    \ext{v}_1
    \EqOrNot
    \left\langle
        \bm a,
        \tensorspan{\ext{\bm y}_a}
    \right\rangle
\) and
\(
    \ext{v}_2
    \EqOrNot
    \left\langle
        \bm b,
        \tensorspan{\ext{\bm y}_b}
    \right\rangle
\)\;

\Return
\(
    (\mathsf{ok},\ext{v}_1,\ext{v}_2,\ext{\bm y}_a, \ext{\bm y}_b)
\)\;
\end{algorithm}

% !TEX root = ../main.tex
\begin{algorithm}[!t]
\caption{Packed Hadamard-product proof}
\label{alg:hadamard}
\KwIn{
\(N=m\kappa\), where \(m,\kappa\in2^{\mathbb N}\);
\(\ext{c}\in\ExtensionField\);
\(\bm a,\bm b\in\Field^N\);
\(\ext{\bm\xi}
=\ext{\bm r}\kronecker\ext{\bm l}\),
where
\(\ext{\bm r}\in\ExtensionField^\kappa\) and
\(\ext{\bm l}
=\tensorspan{\ext{\bm z}_l}\in\ExtensionField^m\)
}
\KwOut{
Terminal values \(\ext{v}_1,\ext{v}_2\) and MLE points
\(\ext{\bm y}_a,\ext{\bm y}_b
\in\ExtensionField^{\log_2N}\);
acceptance bit \(\mathsf{ok}\)
}

Arrange \(\bm a,\bm b\) columnwise as
\(\bm A\leftarrow\concat{\bm a^{(1)},\ldots,\bm a^{(\kappa)}}\)
and \(\bm B\leftarrow\concat{\bm b^{(1)},\ldots,\bm b^{(\kappa)}}
\in\Field^{m\times\kappa}\)\;

\tcp{Encode the column claims in a quotient identity}
\(\Prover\) sets
\(
    \ext{\bm c}^{\top}
    \leftarrow
    \ext{\bm l}^{\top}
    (\bm A\circ\bm B)
    \in\ExtensionField^\kappa.
\)

Using Definition~\ref{def:lagrange}, \(\Prover\) sets:\\
\Indp
\(    \bm g_A(X)
    \leftarrow
    \sum_{j=1}^{\kappa}
    \bm a^{(j)}\psi_j(X)
\);~
\(
    \bm g_B(X)
    \leftarrow
    \sum_{j=1}^{\kappa}
    \bm b^{(j)}\psi_j(X)
\)\;
\(
    g_C(X)
    \leftarrow
    \sum_{j=1}^{\kappa}
    \ext{c}_j\psi_j(X)
\)\;

Set the residual
\(
    f(X)
    \leftarrow
    \left\langle
        \ext{\bm l},
        \bm g_A(X)\circ\bm g_B(X)
    \right\rangle
    -
    g_C(X)
\)
and its quotient
\(
    h(X)
    \leftarrow
    \frac{f(X)}{\Psi(X)}
    \in\ExtensionField[X].
\)
\Indm

\(\PtoV\):
\(\ext{\bm c}\) and the \(\kappa-1\) coefficients of \(h(X)\)\;

\(\Verifier\) computes
\(
    \mathsf{ok}_0
    \leftarrow
    \left(
        \ext{c}
        \EqOrNot
        \left\langle
            \ext{\bm c},\ext{\bm r}
        \right\rangle
    \right).
\)

\(\VtoP\):
\(
    \ext{x}
    \ChooseRandom
    \ExtensionField\setminus\mathcal H.
\)

\tcp{Evaluate to obtain the shorter inner-product claim}
\(\Prover\) and \(\Verifier\) set
\(
    \ext{\bm\rho}
    \leftarrow
    \bigl(
        \psi_1(\ext{x}),
        \ldots,
        \psi_\kappa(\ext{x})
    \bigr)^\top
    \in\ExtensionField^\kappa.
\)

\(\Prover\) sets
\(
    \ext{\bm a}
    \leftarrow
    \bm g_A(\ext{x})
    =
    \bm A\ext{\bm\rho}
    \in\ExtensionField^m
\);~
\(
    \ext{\bm b}
    \leftarrow
    \bm g_B(\ext{x})
    =
    \bm B\ext{\bm\rho}
    \in\ExtensionField^m
\)\;

\(\Verifier\) sets
\(
    \ext{c}_x
    \leftarrow
    g_C(\ext{x})
    =
    \left\langle
        \ext{\bm c},\ext{\bm\rho}
    \right\rangle
\);~
\(
    \ext{d}
    \leftarrow
    \ext{c}_x
    +
    \Psi(\ext{x})h(\ext{x})
\)\;

\(\Prover\) and \(\Verifier\) apply Fact~\ref{fact:ip} to
\(
    \ext{d}
    \EqOrNot
    \left\langle
        \ext{\bm a},
        \ext{\bm b}\circ\ext{\bm l}
    \right\rangle,
\)
and obtain terminal values
\(\ext{u}_1,\ext{u}_2\), the challenge point
\(
    \ext{\bm y}'
    \in\ExtensionField^{\log_2m},
\)
and the acceptance bit \(\mathsf{ok}_{\mathsf{IP}}\)\;

\tcp{Check the packing weights and aggregate the column claims}
For each \(j\in[\kappa]\), \(\Prover\) sets
\(
    \ext{v}_{a,j}
    \leftarrow
    \left\langle
        \bm a^{(j)},
        \tensorspan{\ext{\bm y}'}
    \right\rangle
\);
\(
    \ext{v}_{b,j}
    \leftarrow
    \left\langle
        \bm b^{(j)},
        \ext{\bm l}
        \circ
        \tensorspan{(\ext{\bm y}')}
    \right\rangle
\)\;

Set
\(
    \ext{\bm v}_a
    \leftarrow
    \bigl(
        \ext{v}_{a,1},\ldots,\ext{v}_{a,\kappa}
    \bigr)^\top
\);~
\(
    \ext{\bm v}_b
    \leftarrow
    \bigl(
        \ext{v}_{b,1},\ldots,\ext{v}_{b,\kappa}
    \bigr)^\top.
\)

\(\PtoV\):
\(\ext{\bm v}_a,\ext{\bm v}_b\)\;

\(\Verifier\) checks:
\(
    \mathsf{ok}_{\mathsf{Lag}}
    \leftarrow
    \left(
        \ext{u}_1
        \EqOrNot
        \left\langle
            \ext{\bm v}_a,\ext{\bm\rho}
        \right\rangle
    \right)
    \wedge
    \left(
        \ext{u}_2
        \EqOrNot
        \left\langle
            \ext{\bm v}_b,\ext{\bm\rho}
        \right\rangle
    \right).
\)

\(\VtoP\):
\(
    \ext{\bm{y}}''
    \ChooseRandom
    \ExtensionField^{\log_2 \kappa}.
\)

\(\Prover\) and \(\Verifier\) compute:
\(
    \ext{v}_1
    \leftarrow
    \left\langle
        \ext{\bm v}_a,
        \tensorspan{\ext{\bm y}''}
    \right\rangle
\);~
\(
    \ext{v}_2
    \leftarrow
    \left\langle
        \ext{\bm v}_b,
        \tensorspan{\ext{\bm y}''}
    \right\rangle
\)\;

\(\Prover\) and \(\Verifier\) set the final MLE points:\\
\Indp
\(
    \ext{\bm y}_a
    \leftarrow
    \ext{\bm y}''
    \mathbin{\|}
    \ext{\bm y}'\in\ExtensionField^{\log_2N}
\);~
\(
    \ext{\bm y}_b
    \leftarrow
    \ext{\bm y}''
    \mathbin{\|}
    \left(
        \ext{\bm z}_l
        \circ
        \ext{\bm y}'
    \right)\in\ExtensionField^{\log_2N}
\)\;
\Indm

Output the terminal claims
\(
    \ext{v}_1
    \EqOrNot
    \left\langle
        \bm a,
        \tensorspan{\ext{\bm y}_a}
    \right\rangle
    \wedge
    \ext{v}_2
    \EqOrNot
    \left\langle
        \bm b,
        \tensorspan{\ext{\bm y}_b}
    \right\rangle.
\)\;

Set
\(
    \mathsf{ok}
    \leftarrow
    \mathsf{ok}_0
    \wedge
    \mathsf{ok}_{\mathsf{IP}}
    \wedge
    \mathsf{ok}_{\mathsf{Lag}}.
\)

\Return
\(
    (
        \mathsf{ok},
        \ext{v}_1,
        \ext{v}_2,
        \ext{\bm y}_a,
        \ext{\bm y}_b
    )
\)\;
\end{algorithm}

% !TEX root = ../main.tex
\begin{algorithm}[!t]
\linespread{0.97}\selectfont
\caption{Packed permutation proof}
\label{alg:permutation}
\KwIn{
\(N=m\kappa\), where \(m,\kappa\in2^{\mathbb N}\);
\(\bm a,\bm b\in\Field^N\)
}
\KwOut{
Terminal values \(\ext{v}_a,\ext{v}_b\) and MLE points
\(\ext{\bm y}_a,\ext{\bm y}_b
\in\ExtensionField^{\log_2N}\);
acceptance bit \(\mathsf{ok}\)
}

\tcp{Reduce multiset equality to equality of grand products}
Fix the commitments to \(\bm a\) and \(\bm b\), then sample
\(\VtoP\):
\(\ext z\ChooseRandom\ExtensionField\)\;

\ForEach{\(\bm d\in\{\bm a,\bm b\}\)}{
    Arrange \(\bm d\) columnwise as
    \(
        \bm D^{(\bm d)}
        \leftarrow
        \concat{
            \bm d^{(1)},\ldots,\bm d^{(\kappa)}
        }
        \in\Field^{m\times\kappa}.
    \)

    \(\Prover\) initializes the column-product tree with
    \(
        \bm C^{(\bm d)}_0(Z)
        \leftarrow
        Z\bm 1_{m\times\kappa}
        +
        \bm D^{(\bm d)}
    \)\;

    \For{\(t=0,\ldots,\log_2\kappa-1\)}{
        \(
            \bm C^{(\bm d)}_{t+1}(Z)
            \leftarrow
            \bm C^{(\bm d)}_{t,L}(Z)
            \circ
            \bm C^{(\bm d)}_{t,R}(Z).
        \)
    }

    \(\Prover\) evaluates the last column layer:
    \(
        \ext{\bm c}^{(\bm d)}_0
        \leftarrow
        \bm C^{(\bm d)}_{\log_2\kappa}(\ext z)
        \in\ExtensionField^m.
    \)

    \For{\(i=0,\ldots,\log_2m-1\)}{
        \(
            \ext{\bm c}^{(\bm d)}_{i+1}
            \leftarrow
            \ext{\bm c}^{(\bm d)}_{i,L}
            \circ
            \ext{\bm c}^{(\bm d)}_{i,R}.
        \)
    }

    Set
    \(
        \Theta_{\bm d}
        \leftarrow
        \ext c^{(\bm d)}_{\log_2m}.
    \)
}

\(\PtoV\):
\(\Theta_{\bm a},\Theta_{\bm b}\)\;

\(\Verifier\) checks
\(
    \mathsf{ok}_{\mathsf{perm}}
    \leftarrow
    \left(
        \Theta_{\bm a}
        \EqOrNot
        \Theta_{\bm b}
    \right).
\)

\ForEach{\(\bm d\in\{\bm a,\bm b\}\)}{
    Set
    \(
        \ext{s}^{(\bm d)}_{\log_2m}
        \leftarrow
        \Theta_{\bm d}
    \)
    and
    \(
        \ext{\bm\xi}^{(\bm d)}_{\log_2m}
        \leftarrow(1).
    \)

    \tcp{Unfold the row-product tree}
    \For{\(i=\log_2m-1,\ldots,0\)}{
        Write the current claim as
        \(
            \ext{s}^{(\bm d)}_{i+1}
            \EqOrNot
            \left\langle
                \ext{\bm c}^{(\bm d)}_{i,L},
                \ext{\bm c}^{(\bm d)}_{i,R}
                \circ
                \ext{\bm\xi}^{(\bm d)}_{i+1}
            \right\rangle.
        \)

        Apply Fact~\ref{fact:ip} and randomly combine the two
        terminal claims to obtain
        \(
            \ext{s}^{(\bm d)}_i
            \EqOrNot
            \left\langle
                \ext{\bm\xi}^{(\bm d)}_i,
                \ext{\bm c}^{(\bm d)}_i
            \right\rangle
        \)
        and the acceptance bit
        \(\mathsf{ok}^{(\bm d)}_{\mathsf{row},i}\), retaining
        the tensor factors of \(\ext{\bm\xi}^{(\bm d)}_i\)\;
    }

    Set
    \(
        \ext{q}^{(\bm d)}_{\log_2\kappa}
        \leftarrow
        \ext{s}^{(\bm d)}_0
    \)
    and
    \(
        \ext{\bm\zeta}^{(\bm d)}_{\log_2\kappa}
        \leftarrow
        \ext{\bm\xi}^{(\bm d)}_0.
    \)

    \tcp{Unfold the column-product tree}
    \For{\(t=\log_2\kappa-1,\ldots,0\)}{
        Write the current claim as
        \(
            \ext{q}^{(\bm d)}_{t+1}
            \EqOrNot
            \left\langle
                \ext{\bm\zeta}^{(\bm d)}_{t+1},
                \operatorname{vec}
                \bigl(
                    \bm C^{(\bm d)}_{t,L}(\ext z)
                \bigr)
                \circ
                \operatorname{vec}
                \bigl(
                    \bm C^{(\bm d)}_{t,R}(\ext z)
                \bigr)
            \right\rangle.
        \)

        Apply the coefficient-packing reduction of
        Algorithm~\ref{alg:hadamard} and use Fact~\ref{fact:ip}
        for its evaluated inner product to obtain
        \(
            \ext{q}^{(\bm d)}_t
            \EqOrNot
            \left\langle
                \ext{\bm\zeta}^{(\bm d)}_t,
                \operatorname{vec}
                \bigl(
                    \bm C^{(\bm d)}_t(\ext z)
                \bigr)
            \right\rangle
        \)
        and the acceptance bit
        \(\mathsf{ok}^{(\bm d)}_{\mathsf{col},t}\), retaining
        the tensor factors of \(\ext{\bm\zeta}^{(\bm d)}_t\)\;
    }

    Write
    \(
        \ext{\bm\zeta}^{(\bm d)}_0
        =
        \tensorspan{\ext{\bm y}_{\bm d}}
    \)
    for
    \(
        \ext{\bm y}_{\bm d}
        \in\ExtensionField^{\log_2N}.
    \)

    \(\Prover\) and \(\Verifier\) remove the public shift:
    \(
        \ext{v}_{\bm d}
        \leftarrow
        \ext{q}^{(\bm d)}_0
        -
        \ext z
        \left\langle
            \ext{\bm\zeta}^{(\bm d)}_0,
            \bm 1_N
        \right\rangle.
    \)

    Output the terminal claim
    \(
        \ext{v}_{\bm d}
        \EqOrNot
        \left\langle
            \bm d,
            \tensorspan{\ext{\bm y}_{\bm d}}
        \right\rangle.
    \)
}

Set
\(
    \ext{v}_a\leftarrow\ext{v}_{\bm a}
\),
\(
    \ext{v}_b\leftarrow\ext{v}_{\bm b}
\),
\(
    \ext{\bm y}_a\leftarrow\ext{\bm y}_{\bm a}
\),
and
\(
    \ext{\bm y}_b\leftarrow\ext{\bm y}_{\bm b}.
\)

Set
\(
    \mathsf{ok}
    \leftarrow
    \mathsf{ok}_{\mathsf{perm}}
    \wedge
    \bigwedge_{\bm d\in\{\bm a,\bm b\}}
    \left(
        \bigwedge_{i=0}^{\log_2m-1}
        \mathsf{ok}^{(\bm d)}_{\mathsf{row},i}
    \right)
    \wedge
    \bigwedge_{\bm d\in\{\bm a,\bm b\}}
    \left(
        \bigwedge_{t=0}^{\log_2\kappa-1}
        \mathsf{ok}^{(\bm d)}_{\mathsf{col},t}
    \right).
\)

\Return
\(
    (
        \mathsf{ok},
        \ext{v}_a,
        \ext{v}_b,
        \ext{\bm y}_a,
        \ext{\bm y}_b
    )
\)\;
\end{algorithm}

% !TEX root = ../main.tex
\begin{algorithm}[!t]
\caption{Packed opening protocol}
\label{alg:packed-pcs}
\KwPubIn{
\(\mathsf{pp}_{\pc},\mathsf{cm}\);
\(\Field\subset\ExtensionField\) with basis \(\ext{\bm\beta}\)
and \(\ext\beta_1=1\);
\(N=m\kappa\), \(\ell=\log_2m\);
\(\ext u\in\ExtensionField\),
\(\ext{\bm\rho}\in\ExtensionField^\kappa\), and
\(\ext{\bm\chi}=\bigotimes_{j=1}^{\ell}
\ext{\bm\chi}^{(j)}\), given by factors
\(\ext{\bm\chi}^{(j)}=
(\ext\chi_0^{(j)},\ext\chi_1^{(j)})^\top\in\ExtensionField^2\)
}
\KwPriIn{
\(\Prover\): \(\bm A\in\Field^{m\times\kappa}\) and
commitment randomness \(\rho_{\pc}\)
}
\KwOut{Acceptance bit \(\mathsf{ok}\)}

\(\Verifier\) returns \(0\) on a malformed message or a failed
check, reduction, or opening\;
\tcp{Authenticate the column evaluations by an extension-field projection}
\(\Prover\) sets
\(\ext{\bm a}\leftarrow\bm A\ext{\bm\beta}\) and decodes
\(\bm\Xi\in\Field^{m\times\kappa}\) from
\(\ext{\bm\chi}=\bm\Xi\ext{\bm\beta}\)\;
\(\Prover\) computes \(\bm V\leftarrow\bm A^\top\bm\Xi\)\;
\(\PtoV\): \(\ext{\bm v}\leftarrow\bm V\ext{\bm\beta}
\in\ExtensionField^\kappa\)\;
\(\Verifier\) checks
\(\ext u\EqOrNot\ext{\bm\rho}^{\top}\ext{\bm v}\)\;
\(\Verifier\) recovers \(\bm V\in\Field^{\kappa\times\kappa}\)
from \(\ext{\bm v}=\bm V\ext{\bm\beta}\)\;
\(\VtoP\): \(\ext{\bm r}\ChooseRandom\ExtensionField^\kappa\)\;
\(\Verifier\) computes
\(\ext c\leftarrow\ext{\bm\beta}^{\top}\bm V\ext{\bm r}\)\;
\(\Prover\) computes the operand
\(\bm\Xi\ext{\bm r}\in\ExtensionField^m\)\;
\(\Prover,\Verifier\) apply Fact~\ref{fact:ip} over
\(\ExtensionField\) to
\(\ext c\EqOrNot\langle\ext{\bm a},\bm\Xi\ext{\bm r}\rangle\),
and obtain terminal values \(\ext w_{\mathsf{com}},\ext w\)
and point \(\ext{\bm y}\in\ExtensionField^\ell\)\;
\tcp{The verifier contracts the public queries directly}
Use the multiplication matrices \(\bm M_{\ext z}\) from
Eq.~\eqref{eq:pcs-multiplication-matrix}, applied by extension-field
scaling of rows and columns without materializing \(\bm M_{\ext z}\)\;
\(\Verifier\) initializes
\(\bm W\leftarrow\bm e_1\bm e_1^\top\in\Field^{\kappa\times\kappa}\),
where \(\bm e_1=(1,0,\ldots,0)^\top\)\;
\For{\(j=\ell,\ldots,1\)}{
    \(\Verifier\) updates
    \(\bm W\leftarrow
       \bm W\bm M_{\ext\chi_0^{(j)}}^\top
       +\bm M_{\ext y_j}\bm W
        \bm M_{\ext\chi_1^{(j)}-\ext\chi_0^{(j)}}^\top\)\;
    \tcp{All occurrences on the right use the previous \(\bm W\), as in Eq.~\eqref{eq:pcs-query-recurrence}}
}
\(\Verifier\) checks
\(\ext w\EqOrNot\ext{\bm\beta}^{\top}\bm W\ext{\bm r}\)\;

\(\mathsf{ok}\leftarrow
\pc.\Open\!\left(
    \mathsf{pp}_{\pc},\mathsf{cm},\ext w_{\mathsf{com}},
    \ext{\bm y}:\ext{\bm a},\rho_{\pc}
\right)\)\;
\Return \(\mathsf{ok}\)\;
\end{algorithm}

% !TEX root = ../main.tex
\begin{algorithm}[!t]
\footnotesize
\caption{PLONK proof with single-instance packing}
\label{alg:single-instance-packing}
\KwPubIn{
Parameters of Theorem~\ref{theorem:single-instance-packing};
\(\bm{\mathsf{io}}\);
packed commitments to the five selectors and
\(\{\bm\tau_{\sigma,j}\}_{j=1}^{\gamma}\)
}
\KwPriIn{
\(\Prover\): \(\bm w_L,\bm w_R,\bm w_O\in\Field^N\);
circuit tables and their opening data
}
\KwOut{Acceptance bit \(\mathsf{ok}\)}

\(\Verifier\) returns \(0\) on a malformed message or a failed
check, reduction, or opening\;
\tcp{Bind the witness before sampling challenges}
\ForEach{\(\ell\in\{L,R,O\}\)}{
    \(\Prover\) arranges \(\bm w_\ell\) as
    \(\bm W_\ell\in\Field^{m\times\kappa}\)\;
    \(\PtoV\): a commitment to
    \(\bm W_\ell\ext{\bm\beta}\) as in Section~\ref{sec:pcs}\;
}
\(\Prover\) sets the virtual vector
\(\bm w_M\leftarrow\bm w_L\circ\bm w_R\)\;

\tcp{Check the projected gate equation}
\(\VtoP\): \(\ext{\bm x}\ChooseRandom
\ExtensionField^{\log_2N}\);
set \(\ext{\bm\xi}\leftarrow\tensorspan{\ext{\bm x}}\)\;
\(\PtoV\): the five scalars
\(\ext c_L,\ext c_R,\ext c_M,\ext c_O,\ext c_C\)
from Eq.~\eqref{eq:plonkWeightedHadamard1} and
Eq.~\eqref{eq:plonkWeightedHadamard2}\;
\(\Verifier\) checks
\(\ext c_L+\ext c_R+\ext c_M+\ext c_O+\ext c_C\EqOrNot0\)\;
\ForEach{\(\ell\in\{L,R,O\}\)}{
    Apply the packed Hadamard reduction of
    Section~\ref{sec:hadamard} to
    \(\ext c_\ell\EqOrNot
      \langle\ext{\bm\xi},\bm q_\ell\circ\bm w_\ell\rangle\);
    verify both terminal claims using Algorithm~\ref{alg:packed-pcs}\;
}
Apply the same reduction to
\(\ext c_M\EqOrNot
  \langle\ext{\bm\xi},\bm q_M\circ\bm w_M\rangle\);
open its \(\bm q_M\) claim using Algorithm~\ref{alg:packed-pcs}\;
Write the \(\bm w_M\) terminal claim as
\(\ext u\EqOrNot
  \langle\ext{\bm\rho}\kronecker\ext{\bm\chi},\bm w_M\rangle\),
retaining the weights of Eq.~\eqref{eq:hadamard-terminal}\;
Check this claim by a packed Hadamard reduction with
operands \(\bm w_L,\bm w_R\), column weights \(\ext{\bm\rho}\),
and tensor row weights \(\ext{\bm\chi}\);
verify its two terminal claims using Algorithm~\ref{alg:packed-pcs}\;
Use Algorithm~\ref{alg:packed-pcs} to verify
\(\ext c_C\EqOrNot\langle\ext{\bm\xi},\bm q_C\rangle\)\;

\tcp{Check copy constraints over the intermediate field}
\(\Prover\) forms the padded assignment \(\bm w\in\Field^L\)
and the virtual operands
\(\ext{\bm a}\leftarrow\ext\theta_d\bm w+\ext{\bm\tau}\),
\(\ext{\bm b}\leftarrow\ext\theta_d\bm w+\ext{\bm\tau}_\sigma\)\;
Apply Algorithm~\ref{alg:permutation} to these operands with
base field \(\mathbb K\), extension \(\ExtensionField\),
length \(L\), and packing width \(\kappa_{\mathbb K}\);
its challenge \(\ext z\) is fresh\;
Expand its two terminal claims using
Eq.~\eqref{eq:together-copy-operands}\;
Verify the six wire claims and \(\gamma\) permuted-digit claims
using Algorithm~\ref{alg:packed-pcs}\;
\(\Verifier\) evaluates the public-I/O contributions and
identity labels directly, using
Proposition~\ref{proposition:position-encoding-evaluation},
and checks the two affine recombinations\;
\Return \(1\)\;
\end{algorithm}

% !TEX root = ../main.tex
\begin{algorithm}[!t]
\footnotesize
\linespread{0.97}\selectfont
\caption{PLONK proof with packing across instances}
\label{alg:multiple-instance-packing}
\KwPubIn{
Parameters of Theorem~\ref{theorem:multiple-instance-packing};
the \(\kappa\) public-I/O vectors;
underlying PCS commitments to the five selectors and
\(\ext{\bm\tau}_\sigma\)
}
\KwPriIn{
\(\Prover\): \(\bm W_L,\bm W_R,\bm W_O\in\Field^{m\times\kappa}\);
circuit tables and their opening data
}
\KwOut{Acceptance bit \(\mathsf{ok}\)}

\(\Verifier\) returns \(0\) on a malformed message or a failed
check, reduction, or opening\;
\tcp{Commit in the fixed basis}
\(\Prover\) sets
\(\ext{\bm w}_\ell\leftarrow\bm W_\ell\ext{\bm\beta}\)
for \(\ell\in\{L,R,O\}\), and
\(\ext{\bm w}_M\leftarrow
(\bm W_L\circ\bm W_R)\ext{\bm\beta}\)\;
\(\PtoV\): three underlying PCS commitments to
\(\ext{\bm w}_L,\ext{\bm w}_R,\ext{\bm w}_O\);
retain \(\ext{\bm w}_M\) as a virtual vector\;

\tcp{Check all gate equations by one sum-check}
\(\VtoP\): \(\ext{\bm x}\ChooseRandom
\ExtensionField^{\log_2m}\)\;
Run degree-three sum-check with initial claim \(0\) for the
\(\tensorspan{\ext{\bm x}}\)-weighted residual in
Eq.~\eqref{eq:together-multiple-gates};
obtain terminal point \(\ext{\bm t}\) and claim \(\ext c\)\;
\(\PtoV\): claimed selector evaluations \(\ext v_\ell\)
for \(\ell\in\{L,R,M,O,C\}\), and wire evaluations
\(\ext u_\ell\) for \(\ell\in\{L,R,M,O\}\), at \(\ext{\bm t}\)\;
\(\Verifier\) checks
\(\ext c\EqOrNot
\langle\tensorspan{\ext{\bm x}},\tensorspan{\ext{\bm t}}\rangle
[\sum_{\ell\in\{L,R,M,O\}}\ext v_\ell\ext u_\ell+
(\sum_{j=1}^{\kappa}\ext\beta_j)\ext v_C]\)\;
Verify the five selector and three committed-wire evaluations
by direct underlying PCS openings\;
\tcp{Authenticate the virtual product by a packed Hadamard reduction}
Apply Section~\ref{sec:hadamard}
to
\(\ext u_M\EqOrNot
\tensorspan{\ext{\bm t}}^\top
(\bm W_L\circ\bm W_R)\ext{\bm\beta}\);
verify its two weighted terminal claims using
Algorithm~\ref{alg:packed-pcs}\;

\tcp{Check copy constraints by two extension-field product trees}
\(\Prover\) forms the padded packed assignment
\(\ext{\bm w}\in\ExtensionField^L\), including the public I/O\;
\(\VtoP\): fresh independent
\(\ext y,\ext z\ChooseRandom\ExtensionField\)\;
\(\Prover\) constructs binary product trees on
\(\ext z\bm1_L+\ext{\bm w}+\ext y\ext{\bm\tau}\)
and
\(\ext z\bm1_L+\ext{\bm w}+\ext y\ext{\bm\tau}_\sigma\)\;
\(\PtoV\): the two roots; \(\Verifier\) checks their equality\;
\ForEach{tree}{
    Set the current claim to its root and the current point
    \(\ext{\bm x}\leftarrow()\), with \(\tensorspan{()}=(1)\)\;
    \For{\(\ell=0,\ldots,\log_2L-1\)}{
        Run degree-three sum-check for the current weighted
        parent-product claim, as in
        Appendix~\ref{appx:proof-multiple-instance};
        obtain \(\ext c,\ext{\bm t}\)\;
        \(\PtoV\): child evaluations
        \(\ext v_L,\ext v_R\) at \(\ext{\bm t}\)\;
        \(\Verifier\) checks
        \(\ext c\EqOrNot
        \langle\tensorspan{\ext{\bm x}},
        \tensorspan{\ext{\bm t}}\rangle\ext v_L\ext v_R\)\;
        \(\VtoP\): fresh \(\ext\eta\ChooseRandom\ExtensionField\)\;
        Set the current claim to
        \((1-\ext\eta)\ext v_L+\ext\eta\ext v_R\)
        and \(\ext{\bm x}\leftarrow\ext\eta\mathbin{\|}
        \ext{\bm t}\)\;
    }
    Retain the terminal input claim and its evaluation point\;
}
Expand each terminal claim into its shift, wire, public-I/O,
and label contributions\;
Verify the six wire claims and the one
\(\ext{\bm\tau}_\sigma\) claim by direct underlying PCS openings\;
\(\Verifier\) evaluates the public I/O and identity labels
directly and checks both recombinations, including
the shift \(\ext z\) and label coefficient \(\ext y\)\;
\Return \(1\)\;
\end{algorithm}


\clearpage

\subsection{Packed Polynomial Commitment Scheme}
\label{appx:pcs}

% !TEX root = main.tex
We restate the ring-switching construction of Diamond and
Posen~\cite{eurocrypt/DiamondP26} in the vector and matrix notation
of this paper. This formulation supports the weighted tensor
queries arising from the reductions of Section~\ref{sec:blocks}.

Let \(N=m\kappa\) and \([\ExtensionField:\Field]=\kappa\).
We commit to the fixed-basis packing \(\bm A\ext{\bm\beta}\)
of a matrix \(\bm A\in\Field^{m\times\kappa}\), and reduce
each query to one MLE opening of this commitment.
The verifier evaluates the induced public query by a recursion
on small base-field matrices.
We then instantiate the underlying PCS from a base-field linear
code using Blaze's code-switching construction.

\subsubsection{Given a PCS over the Extension Field}

\medskip\noindent{\bf Commitment and Opening Relation}

Let \(m,\kappa\) be powers of two, with \(\kappa\geq2\), and write
\(\ell=\log_2m\).
Fix an \(\Field\)-basis
\(\ext{\bm\beta}=(\ext\beta_1,\ldots,\ext\beta_\kappa)^\top\)
of \(\ExtensionField\), with \(\ext\beta_1=1\), and let \(\pc\)
be a multilinear PCS for vectors in \(\ExtensionField^m\).
Extension-field elements are stored in this basis, so packing
and coordinate extraction only reorganize their base-field coordinates.
The schoolbook cost bounds below use a polynomial basis.

\paragraph{Setup and commitment.}
Run \(\pc.\Setup(1^\lambda,m)\) to obtain
\(\mathsf{pp}_{\pc}\), and include the fields, basis, and dimensions
\(N,m,\kappa\) in the public parameters.
Pack each row of the witness matrix into one extension-field element:
\[
    \bm A=\concat{\bm a^{(1)},\ldots,\bm a^{(\kappa)}}
        \in\Field^{m\times\kappa},
    \qquad
    \ext{\bm a}=\bm A\ext{\bm\beta}\in\ExtensionField^m.
\]
The commitment is
\(\mathsf{cm}=\pc.\Com(\mathsf{pp}_{\pc},\ext{\bm a};
\rho_{\pc})\), where \(\rho_{\pc}\) is the underlying commitment
randomness. The prover retains \(\bm A\) and \(\rho_{\pc}\).
Since the basis is fixed, binding to \(\ext{\bm a}\) binds the
prover to a unique matrix \(\bm A\in\Field^{m\times\kappa}\).

\paragraph{Opening relation.}
A query specifies column weights
\(\ext{\bm\rho}\in\ExtensionField^\kappa\), a tensor vector
\(\ext{\bm\chi}\in\ExtensionField^m\), and a claimed value
\(\ext u\in\ExtensionField\). The opening relation is:
\begin{equation}
    \ext u
    =\langle\bm A\ext{\bm\rho},\ext{\bm\chi}\rangle
    =\ext{\bm\rho}^{\top}\bm A^\top\ext{\bm\chi}.
\label{eq:pcs-opening-relation}
\end{equation}
The query is supplied by its two-entry tensor factors:
\[
    \ext{\bm\chi}
    =\bigotimes_{j=1}^{\ell}\ext{\bm\chi}^{(j)},
    \qquad
    \ext{\bm\chi}^{(j)}
    =
    \begin{pmatrix}
        \ext\chi^{(j)}_0\\
        \ext\chi^{(j)}_1
    \end{pmatrix}
    \in\ExtensionField^2.
\]
Ordinary MLE queries use the tensor spans of the corresponding
parts of the evaluation point for both \(\ext{\bm\rho}\) and
\(\ext{\bm\chi}\). Allowing arbitrary two-entry factors also
supports the coordinatewise products of tensor queries arising
in the weighted Hadamard reduction.

\medskip\noindent{\bf Packed Opening Protocol}

The commitment uses the fixed weights \(\ext{\bm\beta}\), whereas
the query uses \(\ext{\bm\rho}\). We first claim the column
evaluations, then use coordinate decomposition and a random
projection to reduce their correctness to an inner product with
the committed vector. The verifier evaluates the resulting public
query locally from its tensor factors.

\paragraph{Reduce the column claims.}
Write the public query and the claimed column values in base-field
coordinates:
\[
    \ext{\bm\chi}=\bm\Xi\ext{\bm\beta},
    \qquad
    \ext{\bm v}=\bm V\ext{\bm\beta},
    \qquad
    \bm\Xi\in\Field^{m\times\kappa},
    \quad
    \bm V\in\Field^{\kappa\times\kappa}.
\]
The honest prover computes \(\bm V=\bm A^\top\bm\Xi\) and sends
\(\ext{\bm v}=\bm V\ext{\bm\beta}\).
The verifier extracts \(\bm V\) from this message and checks
\(\ext u=\ext{\bm\rho}^{\top}\ext{\bm v}\).
The matrix \(\bm\Xi\) is defined implicitly by the public query.

Since \(\bm A\) has base-field entries, the remaining claim is
equivalent to a base-field matrix identity:
\[
    \ext{\bm v}=\bm A^\top\ext{\bm\chi}
    \quad\Longleftrightarrow\quad
    \bm V=\bm A^\top\bm\Xi.
\]
Packing its columns with the fixed basis introduces the committed
vector \(\ext{\bm a}\):
\begin{equation}
    \bm V=\bm A^\top\bm\Xi
    \quad\Longleftrightarrow\quad
    \ext{\bm\beta}^{\top}\bm V
    =\ext{\bm a}^{\top}\bm\Xi.
\label{eq:pcs-encoded-matrix}
\end{equation}
After receiving \(\ext{\bm v}\), the verifier samples
\(\ext{\bm r}\sample\ExtensionField^\kappa\), independently of
the original query weights, and projects
Eq.~\eqref{eq:pcs-encoded-matrix} to one inner-product claim:
\begin{equation}
    \ext c
    =\ext{\bm\beta}^{\top}\bm V\ext{\bm r}
    =\langle\ext{\bm a},\bm\Xi\ext{\bm r}\rangle.
\label{eq:pcs-random-opening}
\end{equation}
The verifier computes \(\ext c\) by interpreting the columns of
\(\bm V\) as extension-field elements and taking their inner
product with \(\ext{\bm r}\).

Apply Fact~\ref{fact:ip} over \(\ExtensionField\) to
Eq.~\eqref{eq:pcs-random-opening}.
Let \(\ext{\bm y}\in\ExtensionField^\ell\) be its terminal point
and set \(\ext{\bm\xi}=\tensorspan{\ext{\bm y}}\).
The prover's terminal values
\(\ext w_{\mathsf{com}},\ext w\in\ExtensionField\) give two claims:
\begin{equation}
    \ext w_{\mathsf{com}}
    =\langle\ext{\bm a},\ext{\bm\xi}\rangle,
    \qquad
    \ext w=\ext{\bm\xi}^{\top}\bm\Xi\ext{\bm r}.
\label{eq:pcs-terminal-claims}
\end{equation}
The first is an MLE opening of the committed vector.
The verifier checks the second by the following matrix recursion.

\paragraph{Evaluate the public query.}
Write \(\ext{\bm\xi}=\bm U\ext{\bm\beta}\), where
\(\bm U\in\Field^{m\times\kappa}\).
It suffices to compute the small matrix:
\begin{equation}
    \bm W=\bm U^\top\bm\Xi\in\Field^{\kappa\times\kappa},
    \qquad
    \ext w=\ext{\bm\beta}^{\top}\bm W\ext{\bm r}.
\label{eq:pcs-query-matrix}
\end{equation}
Both \(\bm U\) and \(\bm\Xi\) are determined by public tensor
factors. Their product can be evaluated without forming either
length-\(m\) coordinate table.

For \(\ext z\in\ExtensionField\), let
\(\bm M_{\ext z}\in\Field^{\kappa\times\kappa}\) be the coordinate
matrix of multiplication by \(\ext z\), characterized by:
\begin{equation}
    \ext{\bm\beta}^{\top}\bm M_{\ext z}
    =\ext z\,\ext{\bm\beta}^{\top}.
\label{eq:pcs-multiplication-matrix}
\end{equation}
Thus left multiplication by \(\bm M_{\ext z}\) scales each
coordinate column by \(\ext z\) in \(\ExtensionField\), and right
multiplication by \(\bm M_{\ext z}^\top\) scales each coordinate row.
In particular, \(\bm M_{1-\ext z}=\bm I-\bm M_{\ext z}\).

To expose the recursion, let \(\bm U_j,\bm\Xi_j\) denote the
coordinate matrices for the suffix factors \(j,\ldots,\ell\):
\[
    \bm U_j\ext{\bm\beta}
    =\bigotimes_{i=j}^{\ell}
        \begin{pmatrix}1-\ext y_i\\ \ext y_i\end{pmatrix},
    \qquad
    \bm\Xi_j\ext{\bm\beta}
    =\bigotimes_{i=j}^{\ell}\ext{\bm\chi}^{(i)}.
\]
Here \(\bm U_1=\bm U\), \(\bm\Xi_1=\bm\Xi\), and
\(\bm U_{\ell+1}=\bm\Xi_{\ell+1}=\bm e_1^\top\), where
\(\bm e_1=(1,0,\ldots,0)^\top\in\Field^\kappa\).
The two possible values of the first suffix bit give the block
decompositions:
\begin{equation}
\begin{aligned}
    \bm U_j
    &=
    \begin{bmatrix}
        \bm U_{j+1}\bm M_{1-\ext y_j}^\top\\
        \bm U_{j+1}\bm M_{\ext y_j}^\top
    \end{bmatrix},\\
    \bm\Xi_j
    &=
    \begin{bmatrix}
        \bm\Xi_{j+1}\bm M_{\ext\chi^{(j)}_0}^\top\\
        \bm\Xi_{j+1}\bm M_{\ext\chi^{(j)}_1}^\top
    \end{bmatrix}.
\end{aligned}
\label{eq:pcs-query-blocks}
\end{equation}
Consequently, the verifier computes
\(\bm W_j=\bm U_j^\top\bm\Xi_j\) by initializing
\(\bm W_{\ell+1}=\bm e_1\bm e_1^\top\) and processing
\(j=\ell,\ldots,1\):
\begin{equation}
\begin{aligned}
    \bm W_j
    &=
    \bm M_{1-\ext y_j}\bm W_{j+1}
        \bm M_{\ext\chi^{(j)}_0}^\top
    +\bm M_{\ext y_j}\bm W_{j+1}
        \bm M_{\ext\chi^{(j)}_1}^\top\\
    &=
    \bm W_{j+1}\bm M_{\ext\chi^{(j)}_0}^\top
    +\bm M_{\ext y_j}\bm W_{j+1}
        \bm M_{\ext\chi^{(j)}_1-\ext\chi^{(j)}_0}^\top.
\end{aligned}
\label{eq:pcs-query-recurrence}
\end{equation}
The output is \(\bm W=\bm W_1\).
The verifier never forms the suffix coordinate matrices.
Each step of the second line uses at most \(3\kappa\)
extension-field multiplications, implemented by row and column
scaling without materializing \(\bm M_{\ext z}\).
The recursion takes \(O(\kappa\ell)\) extension-field
multiplications and \(O(\kappa^2\ell)\) base-field additions and
coordinate moves, using \(O(\kappa^2)\) base-field coordinates
of working memory. It applies in every characteristic.

The verifier packs the columns of \(\bm W\), computes
\(\ext{\bm\beta}^{\top}\bm W\ext{\bm r}\), and checks the second
identity in Eq.~\eqref{eq:pcs-query-matrix}.
This computation is exact and requires no additional prover message.

\paragraph{Terminal opening.}
Finally, the parties authenticate \(\ext w_{\mathsf{com}}\)
at \(\ext{\bm y}\) through the underlying PCS:
\[
    \pc.\Open\!\left(
        \mathsf{pp}_{\pc},\mathsf{cm},
        \ext w_{\mathsf{com}},\ext{\bm y}
        :\ext{\bm a},\rho_{\pc}
    \right).
\]
The verifier accepts if the scalar checks, the inner-product
reduction, and the underlying opening accept.
Algorithm~\ref{alg:packed-pcs} in Appendix~\ref{appx:pseudocode}
summarizes the protocol.

\medskip\noindent{\bf Security and Complexity}

\begin{proposition}[Packed polynomial commitment scheme]
\label{proposition:packed-pcs}
Let \(\pc\) bind the prover to a vector in \(\ExtensionField^m\),
with evaluation-soundness error \(\epsilon_{\mathsf{PCS}}\) for
adaptively chosen MLE claims about that vector.
The packed commitment is binding and the opening protocol is
complete, with evaluation-soundness error at most:
\begin{equation}
    \epsilon_{\mathsf{PCS}}
    +\frac{1+2\ell}{|\ExtensionField|}.
\label{eq:pcs-soundness}
\end{equation}
The opening reduction requires no additional lower bound on
\(|\Field|\).

Excluding the underlying PCS opening, the prover's arithmetic cost is:
\begin{equation}
    O(\kappa N)\,\bb
    +O(N)\,\be
    +O(m+\kappa)\,\ee.
\label{eq:pcs-prover-cost}
\end{equation}
The verifier's arithmetic cost is:
\begin{equation}
    O\bigl(\kappa^2(1+\ell)\bigr)\,\bb
    +O\bigl(\kappa(1+\ell)\bigr)\,\ee.
\label{eq:pcs-verifier-cost}
\end{equation}
The verifier bound includes base-field additions and coordinate
reorganization in the matrix recursion.
When \(m\geq\kappa\), schoolbook arithmetic gives
\(O(\kappa N)\) base-field operations for the prover.
Peak prover memory is \(O(N+\kappa^2)\) base-field coordinates,
or \(O(N)\) when \(m\geq\kappa\), excluding the workspace of the
underlying PCS.

With \(\ext u\) already part of the public statement,
prover-to-verifier communication, excluding the underlying
opening, consists of \(\kappa+2\ell+2\) extension-field elements.
The packed commitment inherits hiding from the underlying
commitment. The opening protocol reveals the column evaluations
\(\ext{\bm v}\), so this inheritance does not establish
zero knowledge of the opening transcript.
\end{proposition}

The proof is deferred to Appendix~\ref{appx:proof-pcs}.

\subsubsection{Given a Linear Code over the Base Field}

A linear code over \(\Field\) also encodes packed messages over
\(\ExtensionField\): applying its generator matrix commutes with
basis packing. We use this observation to instantiate the preceding
construction via Blaze's code-switching
protocol~\cite{eurocrypt/BrehmCFRRZ25}, with a BaseFold
code~\cite{crypto/ZeilbergerCF24} as the target. We use Blaze under
its source-encoding and proximity assumptions.

Let \(\Code:\Field^m\to\Field^n\) be a linear code with
generator matrix \(\bm G\in\Field^{n\times m}\).
Its \(\kappa\)-way interleaving encodes the columns of
\(\bm A\in\Field^{m\times\kappa}\) independently.
Using the same fixed basis,
the resulting codeword has the equivalent packed representation:
\begin{equation}
\begin{aligned}
    \bm C&=\bm G\bm A\in\Field^{n\times\kappa},\\
    \ext{\bm c}&=\bm C\ext{\bm\beta}
       =\bm G(\bm A\ext{\bm\beta})
       =\Code_{\ExtensionField}(\ext{\bm a}).
\end{aligned}
\label{eq:pcs-code-packing}
\end{equation}
Here \(\Code_{\ExtensionField}\) is the scalar extension defined
by the same generator matrix; it has the same rate and relative
distance as \(\Code\).
Thus a row of \(\bm C\) represents one extension-field codeword
symbol. Committing to these rows supplies Blaze's source codeword
oracle for \(\ext{\bm a}=\bm A\ext{\bm\beta}\).

Apply Blaze's code-switching construction over \(\ExtensionField\)
to this oracle, switching to a BaseFold encoding of the same
length-\(m\) message \(\ext{\bm a}\).
The resulting MLE openings instantiate the underlying PCS \(\pc\)
in the preceding construction.
The packed opening reduction handles the challenge-dependent weights
\(\ext{\bm\rho}\), and code switching authenticates its terminal
fixed-basis opening. We include the costs and soundness error of
the chosen Blaze instantiation, including source-encoding
consistency, in the underlying PCS costs and
\(\epsilon_{\mathsf{PCS}}\), respectively.

If \(\Code\) is constant-rate and linear-time encodable, forming
\(\bm C\) costs \(O(N)\) base-field operations and its payload
occupies \(O(N\log_2|\Field|)\) bits, excluding Merkle-tree hashes.
The BaseFold stage works on \(m=N/\kappa\) extension-field message
symbols. Thus the source encoding retains the witness's small-field
representation, and the extension-field opening operates on the
shorter packed message.


\clearpage

\subsection{Binary Hadamard Packing in a Mixed Tensor Basis}
\label{appx:binary-hadamard}

Let \(\Field=\mathbb F_2\), \(\kappa=2^s\), and
\(N=m\kappa\), with \([\ExtensionField:\Field]=\kappa\).
We give a direct construction of the \(s\) quotient polynomials
in Eq.~\eqref{eq:overview-multivariate-hadamard} for the Boolean
grid. Unlike the intermediate-field variant of
Section~\ref{sec:hadamard}, this construction keeps the row
polynomials over \(\mathbb F_2\).
The construction retains the quotients in a mixed tensor
basis and uses no general extension-field multiplication.
The row weights \(\ext{\bm l}\in\ExtensionField^m\) are assumed
to have been computed already.

Order the columns of \(\bm A,\bm B\in\Field^{m\times\kappa}\)
lexicographically by \(\{0,1\}^s\), and use the Boolean
Lagrange basis:
\[
    \bm\psi(\bm X)=
       \bigotimes_{\ell=1}^s
       \begin{pmatrix}1-X_\ell\\X_\ell\end{pmatrix},
    \qquad
    \bm g_A(\bm X)=\bm A\bm\psi(\bm X),
    \quad
    \bm g_B(\bm X)=\bm B\bm\psi(\bm X).
\]
For the claimed column products \(\ext{\bm c}\), set
\(g_C(\bm X)=\ext{\bm c}^{\top}\bm\psi(\bm X)\).
The weighted residual and the required certificate are:
\begin{equation}
\begin{aligned}
    f(\bm X)
      &\coloneqq\ext{\bm l}^{\top}
          \bigl(\bm g_A(\bm X)\circ\bm g_B(\bm X)\bigr)
          -g_C(\bm X),\\
    f(\bm X)&=\sum_{j=1}^s(X_j^2-X_j)h_j(\bm X).
\end{aligned}
\label{eq:binary-hadamard-certificate}
\end{equation}
Indeed, \(I(\{0,1\}^s)=
\langle X_1^2-X_1,\ldots,X_s^2-X_s\rangle\).
Thus there are \(s\) quotients rather than a single univariate
quotient.

\smallskip\noindent\textbf{A three-term product identity.}
For \(A(X)=(1-X)a_0+Xa_1\) and
\(B(X)=(1-X)b_0+Xb_1\), the following identity holds over any
characteristic-two coefficient ring:
\begin{equation}
\begin{aligned}
    A(X)B(X)
      &=(1-X)a_0b_0+Xa_1b_1\\
      &\quad +(X^2-X)(a_0+a_1)(b_0+b_1).
\end{aligned}
\label{eq:binary-three-term-product}
\end{equation}
Define the quadratic basis and the binary transform matrix:
\[
    \begin{pmatrix}\chi_0(X)\\\chi_1(X)\\\chi_2(X)\end{pmatrix}
       =\begin{pmatrix}1-X\\X\\X^2-X\end{pmatrix},
    \qquad
    \bm T=\begin{pmatrix}1&0\\0&1\\1&1\end{pmatrix}.
\]
Eq.~\eqref{eq:binary-three-term-product} says that the
coefficient vector of \(A B\) in this basis is
\(\bm T(a_0,a_1)^\top\circ\bm T(b_0,b_1)^\top\).
Tensorizing over all variables gives:
\begin{equation}
\begin{aligned}
    \bm U_A&=\bm A(\bm T^{\otimes s})^{\top},
    &\bm U_B&=\bm B(\bm T^{\otimes s})^{\top}
        \in\Field^{m\times3^s},\\
    \ext{\bm H}&=(\bm U_A\circ\bm U_B)^{\top}\ext{\bm l},
    &\bm\chi(\bm X)&=\bigotimes_{\ell=1}^s
        \begin{pmatrix}
          \chi_0(X_\ell)\\\chi_1(X_\ell)\\\chi_2(X_\ell)
        \end{pmatrix},\\
    \ext{\bm l}^{\top}
        \bigl(\bm g_A(\bm X)\circ\bm g_B(\bm X)\bigr)
      &=\ext{\bm H}^{\top}\bm\chi(\bm X).
\end{aligned}
\label{eq:binary-tensor-product}
\end{equation}
The entries of \(\bm U_A\) and \(\bm U_B\) are bits.
Their coordinatewise products require only AND, and weighting
each product by \(\ext l_i\) means either adding \(\ext l_i\)
or adding zero. These are coefficient transforms, not evaluations
at a third Boolean-field point.

\smallskip\noindent\textbf{Quotient coefficients.}
Index \(\ext{\bm H}\) by
\(\bm t\in\{0,1,2\}^s\).
The \(2^s\) entries with \(\bm t\in\{0,1\}^s\) are exactly
the weighted column products on the Boolean grid. In an honest
execution, they equal \(\ext{\bm c}\) and cancel with \(g_C\).
Assign every remaining entry to the first coordinate \(j\) for
which \(t_j=2\). This defines all quotients directly:
\begin{equation}
    h_j(\bm X)=
    \sum_{\substack{\bm t\in\{0,1,2\}^s\\
                    t_j=2,\ t_1,\ldots,t_{j-1}\in\{0,1\}}}
        \ext H_{\bm t}
        \prod_{\ell<j}\chi_{t_\ell}(X_\ell)
        \prod_{\ell>j}\chi_{t_\ell}(X_\ell).
\label{eq:binary-mixed-quotients}
\end{equation}
In particular, \(h_j\) is independent of \(X_j\), uses the
two-term Lagrange basis in earlier variables, and the three-term
quadratic basis in later variables. Its coefficient slots and
total degree satisfy:
\[
    \text{slots in }h_j=2^{j-1}3^{s-j},
    \qquad
    \deg h_j\leq(j-1)+2(s-j)=2s-j-1.
\]
The total number of coefficient slots is:
\begin{equation}
    Q\coloneqq\sum_{j=1}^s2^{j-1}3^{s-j}
      =3^s-2^s=\kappa^{\log_2 3}-\kappa.
\label{eq:binary-quotient-size}
\end{equation}
For example, \(Q=19,65,665,2059\) for
\(\kappa=8,16,64,128\), respectively.
Grouping the terms of Eq.~\eqref{eq:binary-tensor-product} in
this way gives Eq.~\eqref{eq:binary-hadamard-certificate}
without expanding into monomials or performing polynomial division.

\smallskip\noindent\textbf{Construction cost and storage.}
Apply \(\bm T^{\otimes s}\) recursively rather than as a dense
matrix. Splitting an input into two halves requires \(2^{s-1}\)
XORs to form their sum, followed by three transforms on
\(s-1\) variables. The number \(T_s\) of bit XORs per input row is:
\[
    T_0=0,
    \qquad T_s=3T_{s-1}+2^{s-1},
    \qquad T_s=3^s-2^s=Q.
\]
For each row, compute the transformed bits, AND the corresponding
entries needed by the quotients, and XOR \(\ext l_i\) into each
accumulator whose bit is one. Process rows sequentially; the
matrices \(\bm U_A,\bm U_B\) need not be stored in full.
The counts for generating all quotient coefficients are:
\begin{center}
\begin{tabular}{@{}l@{\quad}r@{\quad}l@{}}
\toprule
Step & Count & Unit \\
\midrule
Transform both inputs & \(2mQ\) & bit XORs \\
Form quotient product coefficients & \(mQ\) & bit ANDs \\
Accumulate weighted coefficients & at most \(mQ\) & extension-field XORs \\
Store all quotients & \(Q\) & extension-field elements \\
\bottomrule
\end{tabular}
\end{center}
Quotient construction uses \(O(m\kappa^{\log_2 3})\) bit
operations. The extension-field XOR count has the same bound;
no general extension-field multiplication is needed.
Equivalently, both counts are
\(O(N\kappa^{\log_2 3-1})\).
If the column products \(\ext{\bm c}\) have not already been
computed, they require another \(m\kappa\) bit ANDs and at most
\(m\kappa\) extension-field XORs, which fit the same asymptotic
bounds. Besides the inputs, supplied weights, and column-claim
vector, a row-streamed implementation stores \(Q\)
extension-field accumulators and
\(O(3^s)\) bits of transform workspace.

An extension-field XOR is not a single-bit operation.
Since an element of \(\ExtensionField\) has \(\kappa\) binary
coordinates, the quotient-construction bound in single-bit units is:
\[
    \text{XORs}\leq(\kappa+2)mQ,
    \qquad
    \text{ANDs}=mQ.
\]
With \(w\)-bit words, accumulation uses at most
\(mQ\lceil\kappa/w\rceil\) word XORs.
These counts exclude loads, selection masks, and writes.
Retaining the mixed tensor representation is part of the bound;
conversion to monomial coefficients would require a separate cost.

\smallskip\noindent\textbf{Reduction after the quotients are fixed.}
The verifier keeps the column-claim check
\(\ext c=\ext{\bm c}^{\top}\ext{\bm r}\) from
Section~\ref{sec:hadamard}.
The prover sends \(\ext{\bm c}\) and all \(Q\) quotient
coefficients before the verifier samples independent uniform
coordinates \(\ext{\bm x}\in\ExtensionField^s\).
The allowed mixed bases enforce the degree bounds above.
Evaluating Eq.~\eqref{eq:binary-hadamard-certificate} yields:
\[
    \left\langle\bm g_A(\ext{\bm x}),
      \bm g_B(\ext{\bm x})\circ\ext{\bm l}\right\rangle
      =g_C(\ext{\bm x})+
       \sum_{j=1}^s(\ext x_j^2-\ext x_j)h_j(\ext{\bm x}).
\]
A false column claim leaves a nonzero residual on the Boolean
grid, regardless of the quotient coefficients. Its total degree
is at most \(2s\), so the identity test has error at most
\(2s/|\ExtensionField|\).
The resulting length-\(m\) inner product and terminal openings
use the existing interfaces. The costs above concern quotient
generation only: generating the row weights, evaluating at the
challenge, and executing the subsequent proof and openings are
separate stages and may require extension-field multiplications.

\begin{proposition}[Binary Hadamard proof: total complexity]
\label{proposition:binary-hadamard-complexity}
Let \(\Field=\mathbb F_2\), \(\kappa=2^s\geq2\), and
\(N=m\kappa\), where \(m\) is a power of two.
Let \(\ExtensionField/\Field\) have degree \(\kappa\), represented
in a polynomial basis with schoolbook arithmetic.
For committed matrices
\(\bm A,\bm B,\bm C\in\Field^{m\times\kappa}\), the mixed-tensor
construction above, followed by Fact~\ref{fact:ip} and three
packed openings from Proposition~\ref{proposition:packed-pcs},
gives a complete proof of \(\bm A\circ\bm B=\bm C\).
Including weight generation and all packed-opening reductions,
but excluding the underlying PCS commitments and openings, its
costs are:
\[
\begin{aligned}
    \text{Prover time}
      &:\quad O\bigl(N\kappa^{\log_2 3}+\kappa^3\bigr)
          \text{ bit operations},\\
    \text{Verifier time}
      &:\quad O\bigl(\kappa^{2+\log_2 3}
              +\kappa^3(1+\log_2m)\bigr)
          \text{ bit operations},\\
    \text{Communication}
      &:\quad O\bigl(\kappa^{\log_2 3}+\log_2m\bigr)
          \text{ extension-field elements},\\
    \text{Prover RAM}
      &:\quad O\bigl(N+\kappa^{1+\log_2 3}\bigr)
          \text{ bits}.
\end{aligned}
\]
If the underlying PCS has evaluation-soundness error
\(\epsilon_{\mathsf{PCS}}\) for adaptively chosen claims, the
soundness error is at most:
\[
    3\epsilon_{\mathsf{PCS}}
    +\frac{3\log_2N+6\log_2m+3}{|\ExtensionField|}.
\]
The excluded costs are those of three underlying commitments and
three underlying openings, including their workspace.
In particular, the prover time is
\(O(N\kappa^{\log_2 3})\) when \(m\geq\kappa\), and the
prover RAM is \(O(N)\) bits when \(m\geq\kappa^{\log_2 3}\).
\end{proposition}

We prove Proposition~\ref{proposition:binary-hadamard-complexity}
in Appendix~\ref{appx:proof-binary-hadamard}.

\clearpage

\subsection{Complexity Bounds for the Hadamard Comparison}
\label{appx:hadamard-comparison}

Table~\ref{tab:hadamard-comparison} compares proofs of
\(\bm a\circ\bm b=\bm c\) for length-\(N\) vectors.
The constraint \(a_i b_i-c_i\) has degree two. Costs include equality
weights, the linear output check, witness-dependent preprocessing, and all reduction
rounds, but exclude commitments and openings. The derived upper bounds
use schoolbook arithmetic: multiplication in
\(\Field\times\Field\), \(\Field\times\ExtensionField\), and
\(\ExtensionField\times\ExtensionField\) costs \(O(1)\),
\(O(\kappa)\), and \(O(\kappa^2)\), respectively, where
\(\kappa=[\ExtensionField:\Field]\). RAM counts base-field elements
including inputs and read-only tables, not measured RSS.
Public field and basis preprocessing is separate; lookup-table
construction is accounted for below. Table accesses are charged
per base-field coordinate, as in Section~\ref{sec:notation}.

\smallskip\noindent\textbf{Optimized rounds and omitted terms.}
Let \(0\leq t\leq\log_2N\) count optimized or skipped rounds.
For small-field-first~\cite{iacr/BagadDT24,iacr/BagadDDT25,iacr/DaoDBDT26},
streamed prefix binding and a weighted \(3^t\)-entry quadratic grid give:
\begin{equation}
\begin{aligned}
\text{Arithmetic:}\quad&
O\!\left(\kappa N(3/2)^t+\kappa^2N/2^t+\kappa^2 3^t\right),\\
\text{RAM:}\quad& O\!\left(N+\kappa N/2^t+\kappa 3^t\right).
\end{aligned}
\label{eq:hadamard-sff-full-cost}
\end{equation}
The last terms account for grid contractions and storage.
For \(N\geq6^t\), they are absorbed by the corresponding
\(N/2^t\) terms, yielding the table entries.
For univariate skip~\cite{iacr/Gruen24}, including FFT interpolation,
weighting, binding, first-round processing, and the sumcheck tail gives:
\begin{equation}
\begin{aligned}
\text{Arithmetic:}\quad&
O\!\left(Nt+\kappa N+\kappa^2N/2^t+\kappa^2 2^t\right),\\
\text{RAM:}\quad& O\!\left(N+\kappa N/2^t+\kappa 2^t\right).
\end{aligned}
\label{eq:hadamard-skip-full-cost}
\end{equation}
Here \(N\geq4^t\) suffices to absorb the last terms.
The required interpolation domains include an order-\(2^t\) subgroup
and sufficiently many off-domain FFT points. When supported,
\(2^t=\Theta(\kappa)\) and \(N=\Omega(\kappa^2)\) give
\(O(\kappa N)\) arithmetic and \(O(N)\) RAM.
For \(t\) fixed by \(\kappa\), the omitted terms are independent of
\(N\). If the absorption conditions fail, use the full bounds in
Eq.~\eqref{eq:hadamard-sff-full-cost} and
Eq.~\eqref{eq:hadamard-skip-full-cost}.

\smallskip\noindent\textbf{Flock's lookup-based skip.}
We re-cost the degree-two zerocheck of
Flock~\cite[Sections~4.2--4.4]{iacr/BunzRW26} with a conventional
sumcheck tail. Fix its byte-sized chunks and subfield
\(\mathbb F_{256}\subseteq\ExtensionField\), and let
\(3\leq t\leq7\) be the number of skipped variables.
Each row has \(2^t\) input bits and \(2^t\) off-domain samples.
Its LDE uses \(2^t/8\) lookups of length-\(2^t\) byte vectors,
so all rows cost \(O(N2^t)\) bit operations, including reads
and XORs. Byte-field products cost \(O(N)\).
In compatible \(\ExtensionField/\mathbb F_{256}\) coordinates,
weighted aggregation costs \(O(\kappa N)\), as do binding at
the skip challenge and the linear output check.
The conventional tail and its equality weights cost
\(O(\kappa^2N/2^t)\). Including parameter-only message
processing and fixed lookup tables gives:
\begin{equation}
\begin{aligned}
\text{Arithmetic:}\quad&
O\!\left(N2^t+\kappa N+\kappa^2N/2^t+\kappa^2 2^t\right),\\
\text{RAM:}\quad&
O\!\left(N+\kappa N/2^t+\kappa 2^t\right).
\end{aligned}
\label{eq:hadamard-flock-full-cost}
\end{equation}
The last terms are absorbed when \(N\geq4^t\).
This bound covers Flock's algebraic reduction to evaluation claims.
Its implementation uses \(t=6\), \(\kappa=128\),
a 16\,KiB compressed LDE table, and a 64\,KiB table for its
fixed equality weights. Compression improves table storage and
cache locality, not the asymptotic LDE operation count.
The admissible skip length is bounded by the byte-field domain,
independently of \(\kappa\).

\smallskip\noindent\textbf{Binary-tower instantiations.}
Assume \(N\geq\kappa\). The Dao--Thaler row pairs its reduction
with conventional sumcheck~\cite{iacr/DaoT24}.
The Liu--Zhang row uses every other level of a Wiedemann
tower~\cite{iacr/LiuZ25}, with mixed subfield/top-field multiplication
and streamed deterministic weights~\cite{iacr/DaoT24}.
RAM retains folded inputs and random-suffix weight tables.

\smallskip\noindent\textbf{Packing assumptions.}
The RMFE--Aurora bound~\cite{iacr/CascudoG21} includes sparse Boolean
R1CS and Reed--Solomon encoding. It assumes constant-rate packing and
codes, linear-space quasilinear FFTs, security \(O(\kappa)\), and a
packed instance large enough to absorb security padding.
Packed sumcheck~\cite{iacr/WeiWXZDWZ25} uses total packed length \(N\),
packing/repetition parameter \(r\), and \(|\Field|>3(2r-1)\).
Streaming forward differences retains \(O(r^2)\) first-round data.
This term remains in the table because \(r\) is independent of
\(\kappa\); it is absorbed if \(N\geq r^2\).
For the base-field \methodname{} row, \(|\Field|\geq\kappa\),
\(N=m\kappa\), \(m\geq\kappa\), and precomputed
bases give \(O(\kappa N)\) arithmetic and
\(O(N+\kappa^2)=O(N)\) RAM, including the linear output check.
All parameter choices must satisfy the target soundness bound.
In particular, \(r=\Theta(\kappa)\) is not assumed, and univariate
skip's error depends on \(2^t\).

\smallskip\noindent\textbf{Intermediate-field accounting.}
Let
\(t_{\mathbb K}=[\mathbb K:\Field]\) be the degree chosen in
Section~\ref{sec:hadamard}. Without lookup tables, multiplication in
\(\Field\times\mathbb K\) costs \(O(t_{\mathbb K})\)
base-field operations, whereas multiplication in \(\mathbb K\)
costs \(O(t_{\mathbb K}^2)\).
In Eq.~\eqref{eq:hadamard-small-field-aggregation}, computing
\(\bm A\bm P\) and \(\bm B\bm P\) therefore costs
\(O(N\kappa t_{\mathbb K})\), and their entrywise product costs
\(O(Nt_{\mathbb K}^2)\).
The contraction with \(\bm L\) uses
\(O(N\kappa/t_{\mathbb K})\) multiplications in \(\mathbb K\),
or \(O(N\kappa t_{\mathbb K})\) base-field operations.
Only the \(\kappa/t_{\mathbb K}\) aggregate coordinate tables
are interpolated. Dense interpolation and division on each cost
\(O(\kappa^2)\) operations in \(\mathbb K\), for a total of
\(O(\kappa^3t_{\mathbb K})\) base-field operations.
The column claims, subsequent inner-product reduction, and linear
output check cost \(O(N\kappa+\kappa^3)\).
Thus Eq.~\eqref{eq:hadamard-small-field-cost} includes the
online parameter-only work; \(N\geq\kappa^2\) gives the
untabulated bound \(O(N\kappa\log_q\kappa)\).

For the space bound, generate one column of \(\bm P\) at a time
using prefix and suffix products of the node differences, and
scan the original rows to accumulate its column of \(\bm U\).
With the public Lagrange denominators precomputed, generating all
columns costs \(O(\kappa^2)\) operations in \(\mathbb K\),
absorbed by the bounds above. The witness and \(\bm L\) use
\(O(N)\) base-field coordinates, and \(\bm U\) uses
\(O(\kappa^2)\). Interpolation basis polynomials can likewise
be generated one at a time from the public vanishing polynomial,
so no dense interpolation matrix is stored.
Fixed conversions between the tower coordinates and the original
commitment basis take at most \(O(N\kappa+\kappa^3)\) operations
and \(O(\kappa^2)\) additional coordinates.
The resulting RAM bound is \(O(N+\kappa^2)\), with the witness
still stored over \(\Field\).

\smallskip\noindent\textbf{Tabulated intermediate-field implementation.}
The second \methodname{} row uses
Eq.~\eqref{eq:hadamard-block-tabulated-cost}.
Choose \(\tau=\lfloor\log_q\kappa\rfloor\), so
\(q^\tau\leq\kappa\) and \(t_{\mathbb K}/\tau=O(1)\).
The block tables of Eq.~\eqref{eq:hadamard-block-table}
form all samples in \(O(N\kappa t_{\mathbb K}/\tau)
=O(N\kappa)\) base-field operations.
Each multiplication-table lookup reads and writes
\(O(t_{\mathbb K})\) coordinates. Thus the
\(O(N\kappa/t_{\mathbb K})\) products in the contraction
cost \(O(N\kappa)\), and aggregate interpolation and division
cost \(O(\kappa^3)\). Samplewise products and the remaining
reduction steps are absorbed in these bounds. Hence:
\begin{equation}
\begin{aligned}
T_{\mathrm{online}}&=O(N\kappa+\kappa^3),\\
T_{\mathrm{pre}}&=O(\kappa^3+t_{\mathbb K}^2q^{2t_{\mathbb K}}),\\
S_{\mathrm{total}}&=O(N+\kappa^3+t_{\mathbb K}q^{2t_{\mathbb K}}).
\end{aligned}
\label{eq:hadamard-tabulation-full-cost}
\end{equation}
The multiplication table has \(q^{2t_{\mathbb K}}\) entries;
schoolbook construction costs \(O(t_{\mathbb K}^2)\) per entry.
The block tables use \(O(\kappa^3)\) coordinates and operations:
precompute scalar multiples of the sample-map rows, then enumerate
input blocks by extending prefixes. Both tables depend only on
public field and interpolation parameters and are reused across
witnesses. Streaming witness rows leaves \(O(N+\kappa^2)\)
working coordinates. For \(N\geq\kappa^2\), the online time is
\(O(N\kappa)\), but the read-only tables remain in the RAM entry.
